\begin{frame}
	\frametitle{time\_steps ($T$) e $\Delta t$: o que são?}
	\framesubtitle{Discretização do tempo contínuo em redes neurais de pulso}
	\small
	\par Redes biológicas e sinais físicos operam em \textbf{tempo contínuo}. Ao simular uma SNN em computadores digitais, o tempo precisa ser fatiado em intervalos discretos --- como os quadros de um filme.
	\begin{center}
		\footnotesize
		\begin{tabular}{lcc}
			\toprule
			& \textbf{time\_steps ($T$)} & \textbf{Passo de tempo ($\Delta t$)} \\
			\midrule
			\textbf{Conceito} & Quantidade de iterações/quadros & Duração física de cada quadro \\
			\textbf{Responde a} & Quantos passos calcular? & Quanto tempo passa por passo? \\
			\textbf{Exemplo} & $T = 16$ passos & $\Delta t = 1{,}0\,\text{ms}$ por passo \\
			\bottomrule
		\end{tabular}
	\end{center}
	\begin{itemize}
		\setlength{\itemsep}{1pt}
		\item \textbf{Janela temporal simulada total}: $t_{\mathrm{total}} = T \cdot \Delta t$.
		\item Com $T=16$ e $\Delta t=1\,\text{ms}$, simula-se $16\,\text{ms}$. Se $\Delta t = 0{,}1\,\text{ms}$, os mesmos 16 passos cobrem apenas $1{,}6\,\text{ms}$ de física real.
		\item $T$ dita o \textbf{esforço computacional} (laços no tempo); $\Delta t$ dita a \textbf{escala temporal física} da dinâmica neuronal \cite{gerstner2014neuronal, eshraghian2023training}.
	\end{itemize}
\end{frame}

\begin{frame}[shrink]
	\frametitle{time\_steps ($T$) e $\Delta t$: qual a sua importância?}
	\framesubtitle{Impacto na memória do neurônio LIF e na estrutura dos tensores}
	\small
	\begin{itemize}
		\setlength{\itemsep}{2pt}
		\item \textbf{Memória e decaimento do potencial ($\beta$)}:
		No neurônio LIF, a fração do potencial que sobrevive entre passos sucessivos é $\beta = \exp(-\Delta t / \tau)$. Com constante de tempo $\tau = 5\,\text{ms}$:
		\begin{itemize}
			\item Se $\Delta t = 1\,\text{ms}$: $\beta = e^{-0{,}2} \approx 0{,}82 \implies$ a membrana retém 82\% da carga (memória temporal ativa).
			\item Se $\Delta t = 16\,\text{ms}$: $\beta = e^{-3{,}2} \approx 0{,}04 \implies$ o potencial vaza quase todo em um único passo (perda de memória).
		\end{itemize}
		\item \textbf{Custo computacional e BPTT}: Mais passos ($T$ alto) capturam dinâmicas mais ricas, mas elevam o custo de tempo e de memória de treino linearmente ($O(T)$).
		\item \textbf{Armadilha na ordem dos tensores}: $(T,B,F)$ guarda a linha $t \cdot B + b$; $(B,T,F)$ guarda a linha $b \cdot T + t$ --- duas ordens diferentes para os MESMOS dados. Exemplo com $T=2, B=2$: a ordem real é $[b_0t_0,\, b_0t_1,\, b_1t_0,\, b_1t_1]$. Se o código espera $(T,B,F)$ mas recebe $(B,T,F)$, a linha 1 --- de verdade $b_0t_1$, o \textbf{segundo} passo da amostra 0 --- é lida como "amostra 1 no passo $t{=}0$": o LIF encadeia a memória de \textbf{duas amostras diferentes} como se fosse uma única linha do tempo. Sem erro, sem aviso --- o treino roda e aprende algo errado (falha silenciosa).
	\end{itemize}
	\par\vfill{\tiny Referências: \cite{gerstner2014neuronal, eshraghian2023training}.}
\end{frame}
